\section{Divisorial types and the cubic flatness defect}
\label{app:cubic-normalization-defect}

This appendix corrects the divisorial sheet classification and records the
part of the quadratic-resolvent method that is currently rigorous.  The former
``smooth cubic-axis theorem'' is replaced by an exact comparison criterion;
the construction and identification of the required maximal
Cohen--Macaulay module remain separate hypotheses.

\subsection{Divisorial types and generic monodromy}

Let \(F\colon X=\A^3\to Y=\A^3\) be a generic-degree-three Keller map, and
let \(\pi\colon\overline X\to Y\) be the finite normalization of \(Y\) in
\(\C(X)\).  Write \(D=\overline X\setminus X\).

\begin{proposition}[Corrected divisorial classification]
\label{prop:cubic-divisorial-trichotomy}
\label{prop:corrected-divisorial-types}
At the geometric generic point of every prime divisor \(H\subset Y\), exactly
one of the following types occurs:
\begin{enumerate}[label=\(U_\arabic*\),start=0]
\item all three sheets are unramified and retained;
\item all three sheets are unramified and exactly one is deleted;
\item all three sheets are unramified and exactly two are deleted;
\end{enumerate}
or
\begin{description}
\item[\(B\)] the inertia is a transposition, the ramified point is deleted,
and the unramified point is retained.
\end{description}
Outside \(S_F\) only \(U_0\) occurs, while at the generic point of a divisor
contained in \(S_F\) only \(U_1,U_2,B\) occur.  In particular, a deleted
three-cycle cannot occur.  The Galois closure of the generic cubic extension
has group \(S_3\).
\end{proposition}

\begin{proof}
The omitted set has codimension at least two by the elementary argument in the
proof of \cref{thm:omitted-singular}.  Hence the geometric generic point of
\(H\) has at least one preimage in \(X\).  Because \(F\) is étale, any
retained branch is unramified.  In the permutation action on the three sheets,
the divisorial inertia group must therefore fix a sheet.  Its only
possibilities are the trivial group and a transposition.

For trivial inertia all three sheets are unramified, and at least one is
retained; this gives \(U_0,U_1,U_2\).  For transposition inertia, the orbit of
size two is ramified and therefore cannot meet \(X\), whereas the fixed sheet
must be retained because the generic point of \(H\) is not omitted.  This is
type \(B\).  A three-cycle has no fixed sheet and is impossible.  Moreover,
\(H\not\subset S_F\) precisely when no sheet is generically deleted and no
ramification occurs, which is type \(U_0\).

The cubic field extension has transitive Galois group, hence its Galois
closure group is either \(C_3\) or \(S_3\).  If it were \(C_3\), every
nontrivial inertia group would be a three-cycle, which was just excluded.
Thus the finite normalization would be unramified in codimension one.  Purity
of the branch locus (Stacks Project, Lemma 58.21.4, Tag
\href{https://stacks.math.columbia.edu/tag/0BMB}{0BMB}) would make it finite
étale over \(\A^3\).  By comparison with the simply connected analytic
space \(\C^3\), affine three-space has no nontrivial connected finite
étale cover.  This contradicts the degree-three field extension.
Therefore the Galois group is \(S_3\).
\end{proof}

\subsection{The quadratic resolvent and an MCM criterion}

Let \((A,\mathfrak m)\) be a complete regular local \(\C\)-algebra of
dimension three, let \(K/\operatorname{Frac}(A)\) be a separable cubic field
extension, and let \(B\) be the integral closure of \(A\) in \(K\), assumed
finite over \(A\).  Let \(\widetilde K\) be the \(S_3\)-Galois closure, let
\(T\) be the integral closure of \(A\) in \(\widetilde K\), and choose a
transposition subgroup \(C_2\subset S_3\) with
\(K=\widetilde K^{C_2}\).  Then
\[
B=T^{C_2}.
\]
Because \(A\) is complete and hence henselian, the finite domains \(B\),
\(T\), and the quadratic resolvent
\[
Q=T^{A_3}
\]
are local.  Assume \(Q\) is normal.  Since \(2\) is invertible and \(A\) is
factorial, the trace-zero summand of the quadratic \(A\)-algebra \(Q\) is
free of rank one; after choosing a generator,
\[
Q\simeq A[w]/(w^2-d)
\]
for some \(d\in A\).  The corrected divisorial classification shows that
\(T/Q\) is unramified in codimension one.  The \(C_3=A_3\)-eigenspace
decomposition therefore has the form
\[
T\simeq Q\oplus L\oplus L^{[2]},
\qquad
L^{[3]}\simeq Q,
\]
where brackets denote reflexive powers.

\begin{proposition}[Resolvent MCM criterion]
\label{prop:resolvent-mcm-criterion}
If \(L\) is maximal Cohen--Macaulay over \(Q\), then the cubic normalization
\(B\) is finite flat over \(A\).
\end{proposition}

\begin{proof}
The hypersurface \(Q\) is Gorenstein.  Since \(L^{[3]}\simeq Q\), one has
\(L^{[2]}\simeq L^\vee\), and the dual of a maximal Cohen--Macaulay module
over a Gorenstein ring is maximal Cohen--Macaulay.  Hence every summand of
\[
T=Q\oplus L\oplus L^{[2]}
\]
is maximal Cohen--Macaulay over \(Q\).  A system of parameters of \(A\) is
also a system of parameters of the finite \(A\)-algebra \(Q\), so it is
\(T\)-regular.  Thus \(\operatorname{depth}_A T=3\).  Auslander--Buchsbaum
over the regular local ring \(A\) gives \(\operatorname{pd}_A T=0\), so
\(T\) is free over \(A\).

Finally, \(B=T^{C_2}\) is an \(A\)-direct summand of \(T\), using the Reynolds
operator \(\frac12(1+\sigma)\).  It is therefore projective, hence free, over
the local ring \(A\).
\end{proof}

\subsection{Three-torsion classes on a threefold hypersurface}

\begin{theorem}[Dao's torsion exclusion and codimension-two detection]
\label{thm:dao-codim-two}
Let \((Q,\mathfrak n)\) be a three-dimensional normal local hypersurface over
\(\C\), and put \(U=\Spec Q\setminus\{\mathfrak n\}\).
\begin{enumerate}[label=(\roman*)]
\item If \(Q\) has an isolated singularity, then \(\Cl(Q)[3]=0\).
\item In general, restriction induces an injection
\[
\Cl(Q)[3]\hookrightarrow
\bigoplus_{\substack{\mathfrak p\in\Sing(Q)\\\operatorname{ht}\mathfrak p=2}}
\Cl(Q_{\mathfrak p})[3].
\]
\end{enumerate}
\end{theorem}

\begin{proof}
Hailong Dao proved that the Picard group of the punctured spectrum of a
three-dimensional local hypersurface is torsion-free; see
\href{https://doi.org/10.1112/S0010437X11005513}{H. Dao, ``Picard groups of
punctured spectra of dimension three local hypersurfaces are torsion-free,''
\emph{Compositio Math.} 148 (2012), 145--152}.

If the singularity is isolated, \(U\) is regular, so rank-one reflexive
modules on \(\Spec Q\) identify with line bundles on \(U\).  Hence
\(\Cl(Q)=\operatorname{Pic}(U)\), and Dao's theorem gives (i).

For (ii), let \([M]\in\Cl(Q)[3]\) restrict trivially at every height-two
singular prime.  At a height-one prime, normality makes \(Q_{\mathfrak p}\) a
DVR; at a regular height-two prime, the local ring is factorial; and at a
singular height-two prime, triviality is the hypothesis.  Thus \(M\) is
locally free at every nonmaximal prime and determines a three-torsion line
bundle on \(U\).  Dao's theorem makes that line bundle trivial.  Reflexive
extension across the closed point then gives \(M\simeq Q\), proving
injectivity.
\end{proof}

\begin{proposition}[Exact comparison criterion]
\label{prop:smooth-axis-comparison}
In the quadratic-resolvent setup above, suppose there exists a rank-one
reflexive \(Q\)-module \(M\) such that
\begin{enumerate}[label=(\roman*)]
\item \(M^{[3]}\simeq Q\);
\item \(M\) is maximal Cohen--Macaulay; and
\item for every height-two singular prime \(\mathfrak p\), the classes of
\(M_{\mathfrak p}\) and \(L_{\mathfrak p}\) agree in
\(\Cl(Q_{\mathfrak p})\).
\end{enumerate}
Then \(L\simeq M\), and the cubic normalization \(B\) is finite flat over
\(A\).
\end{proposition}

\begin{proof}
The class \([L]-[M]\) is three-torsion and restricts to zero at every
height-two singular prime.  By \cref{thm:dao-codim-two}(ii), it is zero in
\(\Cl(Q)\).  Hence \(L\simeq M\), so \(L\) is maximal Cohen--Macaulay.
Apply \cref{prop:resolvent-mcm-criterion}.
\end{proof}

\begin{remark}[Status of the smooth cubic-axis proposal]
A smooth one-parameter family of plane cubics, together with a relative
three-torsion section, suggests a Moore or section module \(M\) satisfying
\(M^{[3]}\simeq Q\).  That geometric picture alone does not prove
\cref{prop:smooth-axis-comparison}(ii)--(iii).  One must still construct an
actual matrix factorization or another depth-three certificate for \(M\), and
then compare its codimension-two local class with the eigensheaf \(L\) coming
from the normalization.  Formal removal of higher transverse terms does not
supply either step.  Accordingly, no unconditional theorem asserting that a
smooth central cubic by itself forces flatness is claimed here.
\end{remark}
